%& -shell-escape
% !TeX root = thesis-text.tex
% !TeX encoding = UTF-8
% !TeX program = XeLaTeX

\documentclass{.local/lib/classes/ndvr-super-article-standalone}

\begin{document}

    % В~течение прошлого десятилетия наблюдался стремительный рост
    % количества видео материалов размещаемых в~Интернете.
    % Вместе с~тем увеличивалось и~вовлечение пользователей
    % в~процесс размещения видео.
    %
    % Значительно увеличилось время пользователя,
    % затраченное на~редактирование, загрузку, поиск и~просмотр видео,
    % Массовая публикация и~обмен информацией породили большое
    % количество нечетких дубликатов видео.


    Начиная с середины 2000-х годов в Интернете
    стало появляться много видео с~очень схожим содержанием видео.
    Такие видео называют почти-дубликатами
    (калька с~английского термина «near duplicate»), полудубликатами,
    неполными или~нечеткими дубликатами.
    Мы будем придерживаться термина «нечеткие дубликаты».
    Слово «нечеткий» в данном случае означает частичное совпадение
    оригинала и дубликата.

    Задача эффективной идентификации нечетких дубликатов играет ключевую
    роль в~задачах поиска, защите авторских прав, и~многих других.

    Последние годы достигнуты серьезные успехи в~поиске нечетких дубликатов,
    выделение низкоуровневых признаков, генерации сигнатур и~индексации видео.

    % В~статье сравниваются существующие варианты
    % определения нечетких дубликатов видео (НДВ),
    % рассматриваются работы по~поиску нечетких дубликатов видео,
    % общие рамки проблемы, современные методы и~появляющиеся тенденции.

\end{document}
